\section{Related Work and Literature Review}
\label{sec:related}

To position LOKA within the current robotics landscape, we analyze related architectures across four primary research domains.

\subsection{Slow-Fast Orchestration and Online Fault Recovery}
\label{sec:related-slowfast}

Decoupling symbolic cognitive reasoning from reactive physical stabilization has been explored across multiple domains.
The AESOP framework utilizes a high-frequency anomaly classifier to trigger a slow generative LLM reasoner, relying on an MPC to preserve physical feasibility during the reasoning latency gap on embedded hardware~\cite{AESOP}.
The REFLECT framework processes multisensory log data into a hierarchical natural language summary to enable an LLM to explain manipulation failures and guide corrective behavior~\cite{REFLECT}.
While structurally related, REFLECT operates as a post-mortem diagnostic tool for general tasks, whereas LOKA's Telemetry Compressor monitors high-frequency proprioceptive stream inputs to execute online adjustments.

Online parameter modification has been explored in wheeled robotics and autonomous driving.
For instance, the Hey Robot! framework uses an LLM to modulate objective cost weights based on human preferences for mobile navigation, though it lacks structural awareness of the physical model~\cite{HeyRobot}.
Similarly, RobotxR1 exposes MPC cost structures and safety margins as online-tunable parameters that a closed-loop reinforcement learning agent alters via language prompts~\cite{RobotxR1}.
However, RobotxR1 treats the underlying kinematic model as an unchangeable prior.
LOKA expands this authority by allowing structural modifications to the solver's body model.
The integration of neuro-symbolic planners, such as Recover, combines explicit rule-based logic with language agents to isolate physical faults~\cite{Recover}.
LOKA avoids the execution gaps of pure symbolic rule bases by mapping outputs directly into a continuous optimal control loop.

Decoupling processing frequencies has also been applied to vehicle routing and monolithic transformers.
AsyncDriver implements an asynchronous multi-threaded system to separate high-latency driving reasoning from real-time path planning~\cite{AsyncDriver}.
In end-to-end foundation policies, Fast-in-Slow (FiS-VLA) builds a unified dual-system policy that embeds high-frequency execution token heads directly inside a larger, low-frequency VLM backbone~\cite{FiS-VLA}.
LOKA shares these multi-rate design patterns but rejects black-box end-to-end actions, ensuring low-level physical safety by anchoring the fast loop to a mathematically transparent MPC solver.

\subsection{Language-to-Rewards and Code Generation}
\label{sec:related-code}

The dominant paradigm for bridging foundation models with optimal control relies on generating executable Python code.
The Eureka and Text2Reward architectures utilize LLMs as automated programmers to search over and debug dense reward functions for reinforcement learning agents~\cite{Eureka, Text2Reward}.
While highly effective, these frameworks operate strictly as heavy offline synthesis engines requiring thousands of environment rollouts.

For online application, the ASPIRE framework distills successful debugging sessions into a reusable text-based skill library using synchronous programmatic synthesis for manipulation tasks~\cite{ASPIRE}.
Similarly, the Language to Rewards (L2R) framework translates natural language text descriptors into functional reward scripts optimized in a loop by a MuJoCo MPC~\cite{Lang2Reward}.
Video2Reward derives reward metrics from target video behaviors for legged robots~\cite{Video2Reward}.
While expressive, these systems are fundamentally bounded by the autoregressive code generation bottleneck, exposing the system to syntax failures and unrecoverable latency during runtime anomalies.
LOKA diverges sharply from this line of work by replacing code generation entirely with a strictly typed parameter and model mutation engine.

\subsection{LLM-Guided Feedback Control Loops}
\label{sec:related-feedback}

A parallel vector of research integrates language models into active feedback control architectures.
The LLMs-guided Adaptive Compensator utilizes an LLM to dynamically output an error-tracking compensator, verifying the loop via stability criteria on soft and humanoid links~\cite{LLMadaptive}.
At the joint gain level, the LLM-Controller framework dynamically modulates nonlinear feedback properties while relying on the underlying baseline tracking architecture to preserve local stability~\cite{LLMcontroller}.
Automated frameworks like ControlAgent and AgenticControl leverage multi-agent language systems to design classical proportional-integral-derivative (PID) layouts by reasoning through phase margins and settling times~\cite{ControlAgent, AgenticControl}.

To bridge optimization with neural safety, architectures such as Never Too Rigid to Reach enforce Lyapunov-constrained reinforcement learning priors to guide virtual model control on flexible arms~\cite{NeverRigid}.
Furthermore, the Meta-Control architecture synthesizes entirely new control topologies depending on changing operational skills~\cite{MetaControl}.
LOKA shares the philosophy of leveraging language priors to modulate control characteristics, but shifts the intervention point.
Rather than operating at the joint-level error feedback or gain layer, LOKA mutates the predictive horizon, task-space priorities, and kinematic assumptions of an online preview optimization engine.

\subsection{Alternative Robot-Language Interface Modalities}
\label{sec:related-interfaces}

Exposing intermediate control spaces to language models remains a highly active area of study.
The SayTap framework utilizes binary foot-contact matrices as a structured intermediate interface to map language commands to quadrupedal locomotion policies~\cite{SayTap}.
OmniVIC couples VLMs with retrieval-augmented generation to continually reconfigure the stiffness and damping properties of a Variable Impedance Controller (VIC), applying explicit force bounds for safety~\cite{OmniVIC}.
For manipulation, the IROSA framework maps user prompts to the operational parameters of Kernelized Movement Primitives (KMPs) through tool calls, completely isolating the language agent from direct trajectory generation~\cite{IROSA}.

To handle complex contact profiles, the CompliantVLA-adaptor and GenCHiP frameworks inject language-derived force compliance, variable impedance values, and compliance boundaries directly into manipulation matrices to ensure safety during physical interaction~\cite{CompliantVLA, GenCHiP}.
LOKA generalizes this mapping philosophy to the full-body locomotion of underactuated bipeds, utilizing a structured YAML configuration block to map semantic objectives directly into high-degree-of-freedom tracking priorities.
